\section{Evaluation}
\label{sec:evaluation}

The evaluation uses the New Delhi platform in
Section~\ref{sec:new_delhi_platform} and asks whether the geometry of
\(\widetilde H_\Phi\) predicts which coefficient, activity, and grouped-source
claims are defensible.  A shared base configuration is changed along one
controlled axis at a time; this avoids conflating wind, source geometry,
background flexibility, noise, and forward-model error.

\subsection{Experiment Matrix}
\label{subsec:experiment_matrix}

\paragraph{Experiment 1: conditioning predicts recovery.}
We vary source geometry and Gaussian observation noise at 0\%, 1\%, 5\%, 10\%,
and 20\% of the maximum clean sensor signal.  For each trial we compare
coefficient and reconstructed-activity errors with \(\sigma_J\), numerical and
effective rank, condition number, and visibility.

\paragraph{Experiment 2: coherent sources require grouped reporting.}
We form increasingly coherent source pairs by shifting a copy of one inventory
map by a controlled spatial offset, so that decreasing the shift drives coherence
toward one.  Individual-source errors are compared with errors in the summed
activity and fitted sensor contribution of each recommended connected component.
Edge reports retain the triggering source--basis pair.

\paragraph{Experiment 3: background correction can help or hurt.}
We compare no background, the normal rank-four basis from
Section~\ref{subsec:new_delhi_forward_models}, a redundant-column basis with the
same span, and a labeled source-like stress basis.  This tests background
flexibility directly: too little background leaves broad confounding in the
residual, while a too-flexible or source-like basis can absorb source-driven
signal and reduce identifiability.  The response geometry is audited through
visibility and the removed/raw absorption ratio.  The rank-four primary and every
sensitivity variant are declared before fitting; no basis is selected from \(Y\)
or from its source-recovery score.

\paragraph{Experiment 4: wind diversity and sensor geometry change resolution.}
Matched source, basis, and sensor sets are evaluated under constant,
single-direction, diurnal, AR(1), multi-directional, and real imputed New Delhi
wind.  Sensor variants are the regulatory layout, seeded random layouts,
downwind-focused layouts, and layouts selected to increase \(\sigma_J\) or
reduce maximum eligible coherence.  The same columns are required on both
sides of every wind comparison.  Separate historical and simulated window
ensembles estimate the 5th, 50th, and 95th percentiles of \(\sigma_J\), rank
probabilities, coefficient weakness probabilities, pairwise ambiguity
probabilities, and conservative component frequencies.

\paragraph{Experiment 5: transport error is amplified by ill-conditioning.}
We use two kinds of transport error.  \emph{Parametric} perturbations vary wind
speed, wind direction, and dispersion within the puff family; a \emph{structural}
mismatch instead generates observations with the auxiliary edge-hold
advection--diffusion PDE while fitting with the open-boundary puff response.  The
operator error norm, coefficient/activity error, residual, and singular spectrum
are reported together, and the parametric transport ensembles may support
transport-uncertainty intervals.  The structural case additionally exercises the
residual-adequacy test of Section~\ref{subsec:model_adequacy}: we report its
rejection rate to show it has power against forward-model misspecification, not
only against missing sources.

\paragraph{Experiment 6: inventory error changes the attribution target.}
We perturb inventory locations, spatial scales, category assignments, and map
versions while holding transport fixed.  Each alternative inventory is a
separate robustness scenario.  Its variation is not pooled with the transport
error of Experiment~5 and is not reported as a confidence interval.

\paragraph{Experiment 7: lag-window sensitivity.}
Physical travel and residence times define the candidate grid.  We evaluate
Equation~\ref{eq:lag_convergence}, choose the smallest lag satisfying
\(\tau_L=10^{-3}\), and report rank, conditioning, and report-component
stability across the grid without using fitted source coefficients for selection.

\paragraph{Experiment 8: missing-source model adequacy.}
Controlled trials omit either a residual-visible source outside
\(\operatorname{span}[H_\Phi^{\mathrm{lag}},Q]\) or an aligned source whose
signature lies in that span.  The refitted parametric-bootstrap test from
Section~\ref{subsec:model_adequacy} is evaluated for calibration and power.  The
aligned case demonstrates why non-rejection cannot certify inventory completeness.

\paragraph{Experiment 9: temporal-basis recovery.}
Known traffic diurnal, brick-kiln intermittent, industry day/night, and mixed
source-basis coefficients are reconstructed at increasing noise.  This study
separates coefficient error from error in the reconstructed source activity
trajectory.

\paragraph{Experiment 10: per-sensor footprints and spatial attribution.}
On controlled trials with known source origins, we check that the per-sensor
footprints of Section~\ref{subsec:per_sensor_footprints} localize the responsible
upwind cells and that per-sensor contributions sum to the fitted sensor signal.
On observed New Delhi data we report, per monitor, the fitted source-group
contributions and their spatial cells of origin, aggregated to the identifiable
report groups so no per-sensor separation exceeds the network's global
resolution.

Real imputed wind is used as a controlled transport input when coefficients are
synthetically known.  The observed-PM\(_{2.5}\) mode is a separate evaluation
and is never assigned synthetic source-recovery metrics.

\subsection{Metrics and Reporting}
\label{subsec:evaluation_metrics}

Controlled trials report absolute and relative coefficient error,
reconstructed-activity error, per-source error, grouped activity and sensor
contribution error, projected sensor residual, uncertainty-interval coverage,
and correlations between diagnostic values and attribution error.  All trials
also report the padded singular spectrum, \(\sigma_J\), numerical and effective
rank, condition number, coefficient visibility, background absorption,
eligible coherence and ray distance, weak flags, triggering coefficient pairs,
and deterministic report groups.  Definitions and undefined weak-pair
conventions are those of Section~\ref{sec:identifiability_theory}.
Lag convergence, the fixed-zero coefficient mask and index mapping, primary
background specification, and ensemble provenance accompany these metrics.
Transport uncertainty and inventory robustness occupy separate fields.

Observed New Delhi runs report raw and projected residuals, fitted sensor
trajectories, normalized proxy coefficients and activities, sensor-space fitted
contributions, observation/transport uncertainty intervals, inventory-scenario
robustness, weak/ambiguous flags, conservative grouped contributions, and the
residual calibration status.  A reported percentage, if used, is explicitly a fraction of the
fitted inventory-attributed sensor signal and not a physical-emission share.

\paragraph{Controlled results.}
The following subsections report the controlled trials, which have known
source--basis coefficients and therefore recovery and coverage metrics.

\subsection{Experiment 1 Results: Conditioning and Recovery}
\label{subsec:results_h1}

\textbf{Results pending.}  The final table will contain noise level, source
geometry, \(\sigma_J\), numerical/effective rank, condition number, coefficient
error, activity error, residual, and interval coverage; the accompanying figure
will plot recovery error against \(\sigma_J\) and condition number.

\subsection{Experiment 2 Results: Coherence and Grouped Reporting}
\label{subsec:results_h2}

\textbf{Results pending.}  The table will report each source pair, triggering
basis pair, eligible coherence, ray distance, individual errors, connected
component, and grouped activity and sensor-contribution errors.  Chain cases
will retain every edge and show when connected components conservatively merge
two sources that are pairwise distinguishable.

\subsection{Experiment 3 Results: Background Stress}
\label{subsec:results_h3}

\textbf{Results pending.}  The table will compare requested/effective background
rank, dependent columns, coefficient visibility, absorption, singular values,
recovery error, and residual for the no-background, normal, redundant, and
source-like stress cases, together with each predeclared basis identifier.

\subsection{Experiment 4 Results: Wind Diversity and Sensor Geometry}
\label{subsec:results_h4}

\textbf{Results pending.}  The table will cross wind provider and sensor layout
with \(\sigma_J\), rank, effective rank, condition number, maximum eligible
coherence, minimum ray distance, weak coefficients, report groups, and activity
error.  Wind comparisons will use identical source--basis columns.  A second
table will report the ensemble quantiles and rank, weakness, ambiguity, and
component probabilities.

\subsection{Experiment 5 Results: Transport Error}
\label{subsec:results_h5a}

\textbf{Results pending.}  The table will identify the true and fitted forward
operators, perturbation magnitude, response error norm, \(\sigma_J\), condition
number, coefficient/activity error, projected residual, and transport-interval
coverage.

\subsection{Experiment 6 Results: Inventory Robustness}
\label{subsec:results_h5b}

\textbf{Results pending.}  The table will identify the primary and alternative
inventory versions, location/scale/category perturbation, diagnostics, fitted
activities, and grouped contributions.  Rows are robustness scenarios, not
confidence-interval draws.

\subsection{Experiment 7 Results: Lag-Window Selection}
\label{subsec:results_lag_selection}

\textbf{Results pending.}  The table will report candidate lag, response
convergence ratio, \(\sigma_J\), numerical/effective rank, condition status,
ambiguity components, selected lag, and the physical candidate-grid rationale.

\subsection{Experiment 8 Results: Missing-Source Adequacy}
\label{subsec:results_missing_source}

\textbf{Results pending.}  The table will report omitted-source alignment,
signal strength, calibrated noise model, \(T_{\mathrm{res}}\), bootstrap
quantile, \(p\)-value, rejection rate, and residual summaries.  It will include
both the residual-visible and span-aligned omissions.

\subsection{Experiment 9 Results: Temporal-Basis Recovery}
\label{subsec:results_temporal_basis}

\textbf{Results pending.}  The table will report source, basis name, true and
estimated coefficients, coefficient error, activity-trajectory error, weak
flags, and uncertainty coverage for traffic, brick kilns, and industries.

\subsection{Experiment 10 Results: Per-Sensor Footprints}
\label{subsec:results_footprints}

\textbf{Results pending.}  On controlled trials the table and plots will report
footprint localization error against known source origins and the per-sensor
contribution reconstruction; the observed New Delhi per-sensor footprints appear
with the apportionment results below.

\paragraph{Observed New Delhi results.}
The remaining subsections report the observed-data study, which has no
source-activity ground truth and therefore reports geometry, residuals,
uncertainty, and conservative groups rather than recovery error.

\subsection{New Delhi Wind-Imputation Validation}
\label{subsec:results_wind_imputation}

\textbf{Results pending.} The table will report held-out station-time vector
error, direction and speed error, mask coverage, query-grid resolution,
checkpoint, seed, device/dtype metadata, and preservation of raw observations in
the audit product.  Because dense real wind truth is unavailable, gridded-field
accuracy is additionally assessed in controlled synthetic wind experiments.

\subsection{New Delhi Identifiability and Report Groups}
\label{subsec:results_new_delhi_identifiability}

\textbf{Results pending.}  The table will report each coefficient's source and
basis, visibility, absorption, weak flag, cross-source trigger pairs, singular
spectrum summary, global warnings, and deterministic connected components.

\subsection{New Delhi Proxy Apportionment and Uncertainty}
\label{subsec:results_new_delhi_apportionment}

\textbf{Results pending.}  The table will report normalized proxy coefficients,
reconstructed activities, fitted sensor-space contributions, active-set and
transport-ensemble intervals, inventory-scenario robustness, and grouped
summaries.  The zero
\texttt{traffic\_06} inventory will be marked unsupported rather than assigned
a substantive estimate.

\subsection{New Delhi Sensor Fit and Residual Diagnostics}
\label{subsec:results_new_delhi_residuals}

\textbf{Results pending.}  The table and plots will report observed-row count,
raw and projected residual norms, station/time residual summaries, fitted versus
observed trajectories, residual autocorrelation, background-removed signal,
response/wind provenance, and---when an external noise model is available---
\(T_{\mathrm{res}}\), bootstrap quantile, \(p\)-value, and adequacy decision.
Without that calibration the result will be labeled \texttt{uncalibrated} and
will not be presented as a model-adequacy pass.
